\documentclass[11pt,a4paper]{article}
\usepackage[T1]{fontenc}
\usepackage{lmodern,amsmath,amssymb,amsthm,mathtools}
\usepackage[margin=27mm]{geometry}
\usepackage{xcolor}
\usepackage[colorlinks=true,linkcolor=blue!45!black,citecolor=blue!45!black,urlcolor=blue!45!black]{hyperref}
\usepackage{microtype}
\newtheorem{theorem}{Theorem}[section]
\newtheorem{proposition}[theorem]{Proposition}
\newtheorem{lemma}[theorem]{Lemma}
\newtheorem{corollary}[theorem]{Corollary}
\theoremstyle{definition}
\newtheorem{definition}[theorem]{Definition}
\newtheorem{remark}[theorem]{Remark}
\newcommand{\OKS}{\mathcal O_{K,S}}
\newcommand{\A}{\mathbb A}
\newcommand{\PP}{\mathbb P}
\setlength{\emergencystretch}{2em}
\hypersetup{pdftitle={Degree and variable lower bounds for Diophantine definitions},pdfsubject={Unconditional obstructions over number fields and rings of S-integers}}
\title{Degree and Variable Lower Bounds\\for Diophantine Definitions}
\author{}
\date{7 November 2026}
\begin{document}
\maketitle
\begin{abstract}
Let $K$ be a number field and let $S$ be a finite set of its finite
primes. A single polynomial whose $K$-rational input projection is
an infinite subset of the ring of $S$-integers must have total
degree at least four and at least three total coordinates. The
coordinate bound also holds for finite systems of arbitrary degree,
even with nonzero conditions and finite alternatives. For a single
polynomial, the input and every essential witness coordinate must
each have individual degree at least two. The proof combines an
elementary argument about prime divisors of polynomial values,
Siegel's theorem for integral points on curves, and Koll\'ar's
unirationality theorem for cubic hypersurfaces. Cubic cones and
geometrically reducible components are treated explicitly.
\end{abstract}

\section{Definitions and statements}

For a nonzero prime ideal $\mathfrak p$ of the ring of integers
$\mathcal O_K$, write $v_{\mathfrak p}$ for its additive valuation,
normalized to have value group $\mathbb Z$, and put
$v_{\mathfrak p}(0)=+\infty$. For a finite set $S$ of nonzero
prime ideals, define
\[
 \OKS=\{a\in K: v_{\mathfrak p}(a)\ge0
                       \text{ for every }\mathfrak p\notin S\}.
\]
In particular, $\mathcal O_{K,\varnothing}=\mathcal O_K$.

\begin{definition}
A single-polynomial presentation over $K$ is an equation
\[
                P(T,Y_1,\ldots,Y_m)=0,
       \qquad P\in K[T,Y_1,\ldots,Y_m].
\]
Its \emph{input projection} is the set
\[
 A_P=\{t\in K:\exists \mathbf y\in K^m, P(t,\mathbf y)=0\}.
\]
The distinguished input $T$ is counted among the coordinates; thus
there are $m+1$ total coordinates and $m$ existential witnesses.
Degree means total degree, including $T$. A witness coordinate is
\emph{essential} if the polynomial depends on it.
\end{definition}

A finite system uses the same coordinates for all its equations.
We also consider finite unions of presentations consisting of equations
and nonzero conditions. Geometrically, each such presentation is a
Zariski open subset of a common zero set. All coordinates are evaluated
in $K$, and no ordering conditions or universal quantifiers are allowed.

\begin{samepage}
\begin{theorem}\label{lb:main}
Let $K$ be a number field and $S$ a finite set of finite primes.
\begin{enumerate}
\item If a single-polynomial presentation has an infinite input
projection contained in $\OKS$, then it has at least three total
coordinates and total degree at least four.
\item The same lower bound of three total coordinates holds for
finite systems of arbitrary degree, with nonzero conditions and finite
alternatives permitted.
\item In the single-polynomial case, the individual degree in $T$
and in each essential witness coordinate is at least two.
\end{enumerate}
\end{theorem}
\end{samepage}

For $K=\mathbb Q$ and $S=\varnothing$, this applies in particular
to every polynomial definition of $\mathbb Z$ over $\mathbb Q$.
The degree restriction in part~(1) concerns one polynomial. A system
of quadratic equations is a different presentation, and is not
excluded by that restriction.

Two established geometric results enter the proof. Siegel's theorem
gives finiteness of integral points on affine curves of positive
genus, including $S$-integral points over a number field
\cite[Chapter~8]{Lang}. Koll\'ar's theorem gives unirationality of a
noncone cubic hypersurface of dimension at least two over a perfect
field as soon as it has a rational point \cite[Theorem~2]{Kollar}.
All number fields are perfect. The other ingredients are proved below.

\section{Rational functions and primes outside a finite set}

\begin{lemma}\label{lb:prime-divisors}
Let $g\in\mathcal O_K[Z]$ be nonconstant. For every finite set
$\Sigma$ of nonzero prime ideals of $\mathcal O_K$, there exist
$a\in\mathcal O_K$ and $\mathfrak p\notin\Sigma$ such that
\[
                  g(a)\ne0,
          \qquad v_{\mathfrak p}(g(a))>0.
\]
\end{lemma}
\begin{proof}
After replacing $g(Z)$ by $g(Z+c)$ for a suitable $c\in\mathcal O_K$,
we may suppose that $b=g(0)\ne0$. Choose a nonzero rational integer
$M$ belonging to every ideal in $\Sigma$; take $M=1$ if $\Sigma$ is
empty. Write $g(Z)=b+c_1Z+\cdots+c_dZ^d$, where $d\ge1$ and $c_d\ne0$.
Then
\[
 h(N)=\frac{g(bMN)}b
      =1+c_1MN+c_2bM^2N^2+\cdots+c_db^{d-1}M^dN^d
      \in\mathcal O_K[N].
\]
For every rational integer $n$ and every $\mathfrak q\in\Sigma$,
we have $h(n)\equiv1\pmod{\mathfrak q}$. The norm polynomial
\[
                    \prod_{\sigma:K\hookrightarrow\mathbb C}h^\sigma(N)
\]
lies in $\mathbb Z[N]$ and has degree $d[K:\mathbb Q]$: its
leading coefficient is the nonzero norm of $c_db^{d-1}M^d$.
Its absolute value is therefore greater than one for all sufficiently
large positive integers $n$. Such an $h(n)$ is nonzero and is not a
unit of $\mathcal O_K$. A prime ideal dividing it lies outside
$\Sigma$ and also divides the nonzero value $g(bMn)=b h(n)$.
Undoing the initial translation gives the required $a$.
\end{proof}

\begin{lemma}\label{lb:rational-function}
Let $r\in K(Z)$ be nonconstant, and let $E\subset K$ be finite.
There is an $a\in K\setminus E$ at which $r$ is defined and
$r(a)\notin\OKS$.
\end{lemma}
\begin{proof}
First suppose that $r$ is a polynomial. Choose a prime ideal
$\mathfrak p\notin S$ at which all its coefficients are integral
and its leading coefficient is a unit. Only finitely many primes
are excluded by these conditions. Choose $0\ne\pi\in\mathfrak p$.
For $a=\pi^{-n}$ and sufficiently large $n$, the leading term of
$r(a)$ has strictly smaller $\mathfrak p$-adic valuation than all
other terms. Its valuation is negative, so
$v_{\mathfrak p}(r(a))<0$. The numbers $\pi^{-n}$ are distinct,
and hence all but finitely many avoid $E$.

Otherwise write $r=f/g$ with relatively prime polynomials
$f,g\in\mathcal O_K[Z]$ and with $g$ nonconstant. The polynomial
Euclidean algorithm over $K$ and clearing denominators give
$A,B\in\mathcal O_K[Z]$ and $0\ne c\in\mathcal O_K$ satisfying
\[
                              Af+Bg=c.
\]
Enlarge $S$ to a finite set $\Sigma$ containing the prime ideals
dividing $c$ or any nonzero value $g(e)$ with
$e\in E\cap\mathcal O_K$.
By Lemma~\ref{lb:prime-divisors}, there are $a\in\mathcal O_K$ and
$\mathfrak p\notin\Sigma$ with $g(a)\ne0$ and
$v_{\mathfrak p}(g(a))>0$. The displayed identity implies
$v_{\mathfrak p}(f(a))=0$. Thus $r(a)$ has negative valuation at
$\mathfrak p\notin S$. Moreover $a\notin E$: if $a=e$ were in
$E\cap\mathcal O_K$, the prime divisor $\mathfrak p$ of its
nonzero value $g(e)$ would have been included in $\Sigma$.
\end{proof}

An integral $K$-variety is called \emph{$K$-unirational} if it is
dominated by an affine space through a rational map defined over $K$.

\begin{proposition}\label{lb:unirational}
Let $X$ be a $K$-unirational integral variety, let $U\subseteq X$
be nonempty and open, and let $f:U\to\A^1_K$ be regular. If
\[
                              f(U(K))\subseteq\OKS,
\]
then $f$ is constant.
\end{proposition}
\begin{proof}
Pull back $f$ through a dominant rational parametrization. If $f$
is nonconstant, the result is a nonconstant rational function on
affine space. Shrink its permitted domain to a nonempty principal
open subset. In characteristic zero, one of its partial derivatives
is nonzero. Since $K$ is infinite, we can fix the other coordinates
in $K$ so that this derivative remains nonzero and none of the
polynomials defining the permitted open set vanishes identically
on the resulting line. The restriction is a nonconstant rational
function $r\in K(Z)$, and the parametrization is permitted at every
$K$-point of the line except a finite set. Lemma~\ref{lb:rational-function}
then supplies a value outside $\OKS$, a contradiction.
\end{proof}

\section{Integral-valued functions on curves}

\begin{lemma}\label{lb:positive-genus}
Let $C$ be a smooth projective geometrically integral curve over
$K$ of positive genus, and let $f\in K(C)$ be nonconstant. There
are only finitely many $P\in C(K)$ at which $f$ is defined and
$f(P)\in\OKS$.
\end{lemma}
\begin{proof}
The function $f$ defines a finite morphism $C\to\PP^1_K$.
Removing its poles gives an affine curve
$D=f^{-1}(\A^1_K)$ finite over $\A^1_K$. Choose generators
$z_1,\ldots,z_r$ of its coordinate ring over $K[f]$. Each $z_i$
satisfies a monic equation over $K[f]$.

There is a finite set $S'$ of finite primes, containing $S$, such
that all coefficients in these finitely many equations belong to
$\mathcal O_{K,S'}[f]$. If $P\in D(K)$ and $f(P)\in\OKS$,
then each $z_i(P)\in K$ is integral over $\mathcal O_{K,S'}$.
That ring is integrally closed, so all $z_i(P)$ belong to it.
Consequently these points are $S'$-integral in a fixed affine
embedding of $D$. Siegel's theorem gives their finiteness
\cite[Chapter~8]{Lang}. Its positive-genus hypothesis holds for
$D$, and at least one pole was removed because $f$ is nonconstant.
\end{proof}

\begin{proposition}\label{lb:curves}
Let $X$ be an algebraic set over $K$ of dimension at most one,
let $U\subseteq X$ be open, and let $f:U\to\A^1_K$ be regular.
If $f(U(K))\subseteq\OKS$, then $f(U(K))$ is finite.
\end{proposition}
\begin{proof}
Pass to the finitely many irreducible components, omitting components
disjoint from $U$. Zero-dimensional components contribute finitely
many points. If an irreducible curve is not geometrically integral,
every $K$-point lies in the intersection of distinct conjugate
geometric components. These intersections form a finite set, so this
case also contributes only finitely many points.

For a geometrically integral curve, take its smooth projective
normalization $C$. Away from finitely many points the normalization
is an isomorphism, and $f$ pulls back to a rational function on $C$.
A constant function has finite image. If the genus is positive,
Lemma~\ref{lb:positive-genus} applies. If the genus is zero and
$C(K)$ is empty, only the finitely many points omitted in passing
to the normalization can contribute. If the genus is zero and
$C(K)$ is nonempty, then $C\simeq\PP^1_K$; Proposition~\ref{lb:unirational}
excludes a nonconstant function with the asserted integral-value
property on the permitted open subset. This proves the result.
\end{proof}

\section{Hypersurfaces of degree at most three}

\begin{theorem}\label{lb:cubic}
Let $X\subseteq\PP^N_K$ be the zero set of a nonzero homogeneous
polynomial of degree at most three. For every open $U\subseteq X$
and regular function $f:U\to\A^1_K$, the containment
$f(U(K))\subseteq\OKS$ implies that $f(U(K))$ is finite.
The same conclusion holds for $X=\PP^N_K$.
\end{theorem}
\begin{proof}
Projective space is unirational, so that case follows from
Proposition~\ref{lb:unirational}. Factor the equation of a hypersurface
over $K$ and treat each irreducible factor separately. Multiplicities
do not change the rational point set.

\emph{Geometrically reducible components.}
An irreducible polynomial of degree two or three over a perfect
field, if not geometrically irreducible, splits geometrically into
conjugate linear factors. Indeed the geometric factors have a common
degree and form one Galois orbit; that degree divides the prime
number two or three. Every $K$-point belongs to all the conjugate
hyperplanes. Their common intersection is a $K$-linear subspace,
and Proposition~\ref{lb:unirational} applies to the restriction of $f$.
Linear components are treated by the same proposition.

\emph{Geometrically integral quadrics.}
A quadric with a smooth $K$-point is rational by projection from
that point. If it has no smooth $K$-point, all its $K$-points lie
in its singular locus, which is a linear subspace. Both cases
give the required finite image.

\emph{Cubics that are not cones.}
Curves and zero-dimensional cases are covered by
Proposition~\ref{lb:curves}. A geometrically integral cubic of
dimension at least two that is not a cone is $K$-unirational if
it has a $K$-point, by Koll\'ar's theorem \cite[Theorem~2]{Kollar}.
This theorem applies to singular cubics as well as smooth ones.
If there is no $K$-point the image is empty; otherwise use
Proposition~\ref{lb:unirational}.

\emph{The structure of a cubic cone.}
Write its cubic equation as $F$ on a $K$-vector space $V$ and put
\[
 W=\{w\in V:\partial_wF=0\text{ as a polynomial}\}.
\]
This is the kernel of a $K$-linear map. In characteristic zero,
$F$ is constant along the directions in $W$, so it descends to
a cubic form on $V/W$. The projective space $\PP(W)$ is the
maximal vertex. Choose a $K$-linear complement to $W$. It gives
a geometrically integral cubic base $B$ which is not a cone.
Indeed, a vertex direction of the base would lift to an additional
direction in the kernel defining $W$. That kernel commutes with
extension of the ground field, so the same assertion holds geometrically.
Its dimension is at least one, since a zero-dimensional cubic
hypersurface is not geometrically integral.

Outside $\PP(W)$, projection from the vertex gives a morphism
to $B$. Each fiber over a $K$-point of $B$ is an affine space of
dimension $\dim W$. On an open subset of $B$ this bundle is
trivial, so the cone is birational to $B\times\A^{\dim W}$.
These statements follow directly by writing a point of $V$ as
the sum of a vector in $W$ and a vector in the chosen complement.

If $B(K)$ is empty, all rational points of the cone lie in the
vertex. If $B(K)$ is finite, apply Proposition~\ref{lb:unirational}
on the permitted open part of each affine-space fiber, and also
on the vertex. This gives a finite union of finite images.
If $\dim B\ge2$ and $B(K)$ is nonempty, Koll\'ar's theorem
makes $B$, and then the cone, unirational.

\emph{A cone over a curve with infinitely many rational points.}
The remaining case has $\dim B=1$ and $B(K)$ infinite. Let $C$
be the smooth projective normalization of $B$. Discarding finitely
many base points, write the function $f$ on an open subset of
$C\times\A^{\dim W}$. If this rational function depends
nontrivially on the affine-space coordinates, some partial
derivative in those coordinates is nonzero. Expressing it as a
quotient of polynomials with coefficients in $K(C)$ shows that,
outside finitely many base points, the specialized function is
nonconstant and its permitted fiber is nonempty. Choose a
$K$-point of $C$ outside those exceptions. Its fiber is a
nonempty open subset of affine space, contradicting
Proposition~\ref{lb:unirational}. Therefore $f$ belongs to $K(C)$.

For every $K$-point of $C$ outside a finite set, the permitted
fiber is a nonempty open subset of affine space and has a $K$-point.
Thus the descended function is defined and $S$-integral at all
these base points. Proposition~\ref{lb:curves} implies that its
image there is finite. Since a nonconstant rational function on
a curve has finite fibers and $C(K)$ is infinite, the function
must be constant. It is then constant as an element of the
cone's function field, and hence wherever it is regular on that
irreducible cone. This includes the points omitted from the
chosen birational chart.
\end{proof}

\section{Consequences for polynomial presentations}

\begin{proof}[Proof of Theorem~\ref{lb:main}(1) and (2)]
First consider at most two total coordinates. For one system of
equations, its zero set is either all of the ambient affine space
or has dimension at most one. In the first case, the permitted
nonempty open subset is unirational, and its input function cannot
have an infinite $S$-integral image by Proposition~\ref{lb:unirational}.
In the second case, Proposition~\ref{lb:curves} gives a finite
input image. Empty permitted sets contribute nothing. Taking
finitely many alternatives preserves finiteness. This proves the
coordinate bound for systems and for single equations.

For a single nonconstant polynomial of total degree at most three,
homogenize it and regard the distinguished input as a regular
function on the usual affine chart of its projective zero set.
Theorem~\ref{lb:cubic} gives a finite input image, independently
of the number of coordinates. A nonzero constant polynomial has
empty projection. The zero polynomial has projection $K$, which
is not contained in $\OKS$ because $S$ is finite. These observations
prove the degree bound.
\end{proof}

\begin{proposition}\label{lb:individual-degrees}
Suppose that $P\in K[T,Y_1,\ldots,Y_m]$ has an infinite input
projection contained in $\OKS$. Then
\[
                \deg_T P\ge2,
    \qquad \deg_{Y_j}P\ge2
          \quad\text{for every essential witness }Y_j.
\]
\end{proposition}
\begin{proof}
Suppose first that $P$ is linear in an essential witness $Y_j$.
Write
\[
             P=A(T,\mathbf Y')Y_j+B(T,\mathbf Y'),\qquad A\ne0.
\]
Choose constants $\mathbf c\in K^{m-1}$ so that
$A(T,\mathbf c)$ is a nonzero polynomial. Such a choice exists
because $K$ is infinite. For every $t$ outside its finite root
set, setting $Y_j=-B(t,\mathbf c)/A(t,\mathbf c)$ gives a
rational witness. The projection is therefore cofinite in $K$.
It cannot be contained in $\OKS$, whose complement is infinite:
for a prime $\mathfrak p\notin S$ and $0\ne\pi\in\mathfrak p$,
all the distinct elements $\pi^{-n}$ lie outside $\OKS$.

Suppose next that $\deg_TP\le1$, and write
$P=A(\mathbf Y)T+B(\mathbf Y)$. If $A=0$, the input projection
is either empty or all of $K$. Otherwise on the nonempty open
set $A\ne0$ the rational function $-B/A$ takes only values in
$\OKS$. Proposition~\ref{lb:unirational} makes it constant.
At any rational point with $A=B=0$, every input would be admitted,
which is impossible. Thus the projection has at most one element,
again a contradiction. This proves Theorem~\ref{lb:main}(3).
\end{proof}

\begin{corollary}\label{lb:linear-systems}
No finite system of linear equations, even with nonzero conditions
and finite alternatives, has an infinite input projection contained
in $\OKS$.
\end{corollary}
\begin{proof}
Each nonempty common zero set of linear equations is an affine
space over $K$. Apply Proposition~\ref{lb:unirational} to its
permitted open subset and take the finite union of the resulting
finite images.
\end{proof}

\begin{corollary}\label{lb:degree-sharpness}
Let $A\subseteq\OKS$ be infinite. If $A$ has a single-polynomial
presentation of degree four over $K$, then four is the least
possible degree of any such presentation, regardless of how many
witness coordinates are allowed. If $A$ has a presentation by a
system of quadratic equations, then two is the least possible
maximum degree of the equations in a system.
\end{corollary}
\begin{proof}
The first assertion follows from Theorem~\ref{lb:main}; the second
follows from Corollary~\ref{lb:linear-systems}.
\end{proof}

The coordinate lower bound permits any number of equations and
arbitrary degrees. In contrast, the degree-four obstruction concerns
the geometry of one hypersurface. Neither statement asserts the
existence of a presentation in three total coordinates.

\begin{thebibliography}{9}
\bibitem{Kollar}
J.~Koll\'ar, Unirationality of cubic hypersurfaces,
\emph{Journal of the Institute of Mathematics of Jussieu}
\textbf{1} (2002), 467--476, Theorem~2.
\url{https://arxiv.org/abs/math/0005146}.

\bibitem{Lang}
S.~Lang, \emph{Fundamentals of Diophantine Geometry},
Springer, 1983, Chapter~8, ``Siegel's Theorem and Integral Points'',
188--224.
\url{https://doi.org/10.1007/978-1-4757-1810-2_8}.
\end{thebibliography}
\end{document}
